% Part: turing-machines
% Chapter: undecidability
% Section: unsolvability-decision-problem

\documentclass[../../../include/open-logic-section]{subfiles}

\begin{document}

\olfileid[fr]{tur}{und}{uns}
\olsection{L'insolubilité du problème de la décision}

\begin{thm}
\ollabel{thm:decision-prob}
Le problème de la décision est insoluble : il n'existe aucune machine
de Turing~$D$ qui, lancée sur un ruban portant en entrée
!!a{sentence} de logique du premier ordre~$!B$, s'arrête toujours
et produit~$1$ si et seulement si $!B$ est valide, et~$0$ sinon.
\end{thm}

\begin{proof}
Supposons que le problème de la décision soit soluble, et soit~$D$
une machine qui le résout. Nous pourrions alors résoudre le problème
de l'arrêt comme suit. Construisons une machine de Turing~$E$ qui,
recevant l'indice~$e$ de la machine~$M_e$ et l'entrée~$w$, calcule
l'!!{sentence} correspondant $!T(M_e,w) \lif !E(M_e,w)$, puis s'arrête,
sa tête lisant la première case du ruban.\footnote{Correction éditoriale : la source fait revenir la tête sur la case la plus à gauche, d'indice $0$. Pour lancer $D$ sur l'énoncé codé, elle doit lire la case $1$, à droite du marqueur. Les alphabets du codeur et du décideur sont rendus communs comme pour le raccordement précédent.} Lancée sur~$e$ et~$w$,
la machine raccordée $E \concat D$ calcule d'abord
$!T(M_e,w) \lif !E(M_e,w)$, puis soumet cet énoncé à~$D$.
Cette dernière s'arrête avec la sortie~$1$ si et seulement si
$!T(M_e,w) \lif !E(M_e,w)$ est valide, et avec la sortie~$0$
sinon. D'après le \olref[ver]{lem:halt-if-valid} et
le \olref[ver]{lem:valid-if-halt}, cet énoncé est valide si et seulement
si~$M_e$ s'arrête sur l'entrée~$w$. Ainsi, $E \concat D$ s'arrête
avec la sortie~$1$ si~$M_e$ s'arrête sur~$w$, et avec la sortie~$0$
sinon : elle résoudrait le problème de l'arrêt. Or, d'après
le \olref[hal]{thm:halting-problem}, aucune machine de Turing ne peut
le résoudre.
\end{proof}

\begin{cor}\ollabel{cor:undecidable-sat}%
Le problème de savoir si un !!{sentence} arbitraire de logique du
premier ordre est satisfaisable est indécidable.
\end{cor}

\begin{proof}
Supposons qu'une machine de Turing~$S$ décide la satisfaisabilité.
Nous pourrions résoudre le problème de la décision de la manière
suivante. Étant donné en entrée !!a{sentence}~$!B$, décalons-le
d'une case vers la droite. Revenons à la case~$1$ et inscrivons-y
le symbole~$\lnot$.

Lançons maintenant~$S$. Elle finit par s'arrêter avec la sortie~$1$
si~$\lnot !B$ est satisfaisable, et avec la sortie~$0$ sinon. Nous choisissons $S$ disciplinée, pour que sa tête lise alors la case~$1$.
Si la case~$1$ porte un symbole~$\TMstroke$, effaçons-le; si elle
est vide, inscrivons-y un symbole~$\TMstroke$, puis arrêtons-nous.

La machine ainsi construite s'arrête toujours. Elle produit~$1$
si et seulement si $\lnot !B$ est insatisfaisable, et~$0$ sinon.
Comme $!B$ est valide si et seulement si $\lnot !B$ est
insatisfaisable, elle produirait~$1$ exactement lorsque $!B$ est
valide, et~$0$ sinon. Elle résoudrait donc le problème de la décision.
\end{proof}

\begin{explain}
Aucune machine de Turing ne peut donc toujours répondre correctement
« oui » ou « non » à la question : « $!B$ est-il un !!{sentence}
valide de logique du premier ordre ? » Il \emph{existe} pourtant une machine
qui répond correctement « oui » lorsque la réponse est affirmative,
mais qui ne s'arrête pas lorsque la réponse est négative. Cela
résulte des théorèmes de correction et de complétude de la logique
du premier ordre, ainsi que de la possibilité d'énumérer
effectivement les !!{derivation}s.
\end{explain}

\begin{thm}
\ollabel{thm:valid-ce}%
La validité des !!{sentence}s du premier ordre est semi-décidable :
il existe une machine de Turing~$E$ qui, lancée sur un ruban portant
en entrée !!a{sentence} de logique du premier ordre~$!B$, s'arrête
avec la sortie~$1$ si et seulement si $!B$ est valide, et ne
s'arrête pas sinon.
\end{thm}

\begin{proof}
Un algorithme effectif peut engendrer successivement toutes les
!!{derivation}s possibles en logique du premier ordre. La
machine~$E$ les énumère et, dès qu'elle trouve !!a{derivation}
établissant $\Proves !B$, elle s'arrête avec la sortie~$1$.
Par le théorème de correction, si~$E$ s'arrête ainsi, alors
$\Entails !B$. Réciproquement, si $\Entails !B$, le théorème de
complétude assure l'existence d'!!a{derivation} établissant
$\Proves !B$. Puisque~$E$ engendre systématiquement toutes les
!!{derivation}s possibles, elle finira par trouver celle-ci et
s'arrêtera avec la sortie~$1$.
\end{proof}

\end{document}
